\documentclass[11pt]{article}
\usepackage{amsmath,amssymb,hyperref}
\usepackage[margin=1in]{geometry}
\title{An exact algebraic three-dimensional second-order anticoherent subspace at seventeen qubits}
\date{29 October 2026}
\begin{document}
\maketitle
\paragraph{Result.} There exists a three-dimensional subspace of the symmetric
seventeen-qubit Hilbert space all of whose normalized vectors are second-order
anticoherent. The coefficients below are specified as an isolated real algebraic
root, with a rational contraction certificate proving existence. This is an exact
existence result, not a conclusion from a small numerical residual. The
auxiliary-spin inertia theorem added on 30 October 2026 gives $4k\le N-2$
for every $N\ge8$. At $N=17$ this yields $k\le3$, so the construction is
globally optimal: $k_{\max}(17,2)=3$. The independent general proof is supplied
in the catalogue's \texttt{research/ancilla-inertia/proof.tex}.

\section{Coefficients and rational defining equations}
Use normalized Dicke states $|D_{17}^{(r)}\rangle$, with $r$ the excitation number,
and put
\[
 s_r=\left\lfloor\sqrt{\binom{17}{r}}+\frac12\right\rfloor.
\]
The ordered supports of three basis vectors are
\[
 R_0=R_1=(2,5,8,11,14,17),\qquad R_2=(0,3,6,9,12,15).
\]
For $a=0,1,2$ and $\ell=0,\ldots,5$, define
\[
 z_{a,R_{a,\ell}}=
 \frac{\xi_{12a+2\ell}+i\xi_{12a+2\ell+1}}{s_{R_{a,\ell}}},
 \qquad
 |u_a\rangle=\sum_{r=0}^{17}\sqrt{\binom{17}{r}}z_{a,r}|D_{17}^{(r)}\rangle,
\]
with $z_{a,r}=0$ outside $R_a$. The coordinates $\xi_i$ are real. Fix
\[
 \xi_1=\xi_3=\xi_{12}=\xi_{13}=\xi_{15}=\xi_{25}=0,
 \qquad
 \xi_9=-\frac{191}{1000},\quad
 \xi_{22}=\frac{41}{500},\quad
 \xi_{32}=-\frac{34}{125}.
\]
The remaining 27 coordinates are the unique root in the rational box specified
in \texttt{certificate.json} of the following 27 real quadratic equations.
Write
\[
 H_{ab}^{ij}=\sum_{\ell=0}^{15}\binom{15}{\ell}
 z_{a,i+\ell}\overline{z_{b,j+\ell}}.
\]
Impose:
\begin{align*}
 H_{aa}^{ii}&=\frac{1}{3\binom2i}
 &&(a=0,1,2;\ i=0,1,2),\\
 H_{01}^{ii}&=0 &&(i=0,1,2),\\
 H_{a2}^{ij}&=0
 &&(a=0,1;\ (i,j)=(0,1),(1,2),(2,0)).
\end{align*}
The first line supplies nine real equations. Real and imaginary parts of the
second and third lines supply six and twelve real equations respectively. All
coefficients are rational because of the chosen Dicke scaling.

The symmetric two-qubit reduction obeys
\[
 \bigl[R_2(|u_a\rangle\langle u_b|)\bigr]_{ij}
 =\sqrt{\binom2i\binom2j}\,H_{ab}^{ij}.
\]
The supports and Hermitian conjugation show that every other required real or
imaginary component is either identically zero or one of the selected equations,
possibly with its sign reversed. The verifier reconstructs all 162 real and
imaginary components and checks this implication exactly. Hence any real root
of the defining system gives
\[
 R_2(|u_a\rangle\langle u_b|)=\delta_{ab}\frac{I_3}{3}.
\]
Taking traces gives $\langle u_b|u_a\rangle=\delta_{ab}$, so the three vectors are
orthonormal and span a genuine three-dimensional subspace. Expanding any
normalized superposition establishes the full-subspace anticoherence condition.

\section{Rational certificate of a root}
Let $F:\mathbb R^{27}\to\mathbb R^{27}$ be the selected polynomial system after
substituting the nine fixed coordinates, in the variable order recorded in the
certificate. The same file supplies an explicit rational center $x_0$, a rational
matrix $A$ approximating $DF(x_0)^{-1}$, and the box radius
\[
 r=10^{-60}.
\]
Every entry is a rational number; the long lists specify the algebraic root
unambiguously and are included in the downloadable certificate. Consider the
map $T(x)=x-AF(x)$. In the infinity norm define
\[
 \eta=\|AF(x_0)\|_\infty,\quad
 \delta=\|I-A\,DF(x_0)\|_\infty,\quad
 H=\max_i\sum_{j,k}\left|\frac{\partial^2F_i}{\partial x_j\partial x_k}\right|,
 \quad q=\delta+r\|A\|_\infty H.
\]
The Hessians are constant because the system is quadratic. Exact rational
arithmetic proves
\[
 \delta<1,\qquad \eta<10^{-90},\qquad q<10^{-56},
 \qquad \eta+qr<r.
\]
Therefore $T$ maps the closed box $\|x-x_0\|_\infty\le r$ into itself and is a
strict contraction there. Banach's theorem gives a unique fixed point.
The inequality $\delta<1$ implies that $A$ is nonsingular, so this fixed point is
a root of $F$. Moreover $\|I-A\,DF(x)\|_\infty<1$ throughout the box, so the root's
Jacobian is nonsingular. An isolated nonsingular solution of a rational
polynomial system has algebraic coordinates. This justifies the exact algebraic
interpretation of the displayed components.

\section{Portable verification and provenance}
Run \texttt{python verify.py} with SymPy available. The verifier uses no numerical
root finder. It reconstructs the rational equations, checks their completeness,
recomputes every contraction bound, and checks that \texttt{basis.json}'s exact
coefficient expressions and decimal previews agree with this same definition.
The function \texttt{verify()} also exposes the result to the catalogue audit.
The \texttt{parameterIntervals} in the basis record include every coordinate,
including the fixed ones, for independent interval checks of derived quantities.

The source's numerical support pattern guided the choice of coordinates and the
initial approximation. Three free coordinates were fixed rationally, so equality
to the original decimal span is not asserted. The exact existence claim rests
on the rational certificate above. No claim of literature priority is made.
\end{document}
